\documentclass[11pt]{article}

\usepackage[letterpaper,margin=1in]{geometry}
\usepackage[T1]{fontenc}
\usepackage[utf8]{inputenc}
\usepackage{lmodern}
\usepackage{microtype}
\usepackage{amsmath,amssymb,mathtools,bm}
\usepackage{xcolor,slashed,needspace}
\usepackage{enumitem}
\usepackage{hyperref}
\usepackage[most]{tcolorbox}
\usepackage{fancyhdr}
\usepackage{parskip}

\definecolor{qftblue}{HTML}{245A8D}
\definecolor{qftgray}{HTML}{F4F6F8}
\hypersetup{
  colorlinks=true,
  linkcolor=qftblue,
  urlcolor=qftblue,
  pdftitle={PHYS 662 Homework 4 Solutions},
  pdfauthor={Nima Lashkari}
}

\pagestyle{fancy}
\fancyhf{}
\lhead{PHYS 662: Quantum Field Theory I}
\rhead{Homework 4 Solutions}
\cfoot{\thepage}
\setlength{\headheight}{14pt}

\setlist[enumerate]{leftmargin=2.2em,itemsep=0.55em,topsep=0.5em}
\newlist{parts}{enumerate}{1}
\setlist[parts]{label=\textbf{(\alph*)},leftmargin=2.3em,itemsep=0.65em}

\newcommand{\dd}{\mathrm{d}}
\newcommand{\ii}{\mathrm{i}}
\newcommand{\R}{\mathbb{R}}
\newcommand{\C}{\mathbb{C}}
\newcommand{\Z}{\mathbb{Z}}
\newcommand{\Lag}{\mathcal L}
\newcommand{\tr}{\operatorname{tr}}
\newcommand{\Ad}{\operatorname{Ad}}

\newtcolorbox{conventionbox}[1][]{
  enhanced,breakable,
  colback=qftgray,colframe=qftblue,
  boxrule=0.8pt,arc=1.5mm,
  title={#1},fonttitle=\bfseries,coltitle=white
}
\newcommand{\problem}[2]{%
  \Needspace{10\baselineskip}
  \section*{\color{qftblue}Problem #1: #2}
  \addcontentsline{toc}{section}{Problem #1: #2}
}
\newcommand{\lectureconnection}[1]{%
  {\small\textit{Lecture connection: #1.}\par}\medskip
}


\newcommand{\diag}{\operatorname{diag}}
\newcommand{\SO}{\mathrm{SO}}
\newcommand{\SU}{\mathrm{SU}}
\newcommand{\SL}{\mathrm{SL}}
\newcommand{\ISO}{\mathrm{ISO}}
\newcommand{\ddPi}{\dd\Pi}
\newcommand{\vac}{\lvert0\rangle}
\newcommand{\vbra}{\langle0\rvert}
\newcommand{\Dc}{\mathcal C}
\newcommand{\one}{\mathbf1}

\usepackage{booktabs,tikz}
\newcommand{\hc}{\mathrm{h.c.}}
\newcommand{\partsol}[2]{\par\Needspace{4\baselineskip}\medskip\noindent\textbf{(#1) #2.}\quad}
\begin{document}
\begin{center}
{\scshape\Large PHYS 662 --- Quantum Field Theory I\par}
\medskip
{\Huge\bfseries\color{qftblue} Homework 4 Solutions\par}
\medskip
{\large Lectures 11--14 \quad Fall 2026\par}
\end{center}
\begin{conventionbox}[Conventions and scope]
These worked solutions cover every part of all eight problems in the
written Homework Set 4, including the two added problems. Use
\(\eta=(+---)\), \(\epsilon^{0123}=1\), \(\epsilon_{123}=1\),
\(J_i=\epsilon_{ijk}M^{jk}/2\), and \(K_i=M^{i0}\).
The mode spinors have ordinary complex entries, while classical field
components in the mass-action problem are Grassmann odd.
\[
 s^\mu=(I,\bm\sigma),\quad \widetilde s^\mu=(I,-\bm\sigma),\quad
 \gamma^\mu=\begin{pmatrix}0&\widetilde s^\mu\\s^\mu&0\end{pmatrix},
 \quad E=i\sigma_2,\quad \gamma_*=-i\gamma^0\gamma^1\gamma^2\gamma^3.
\]
Signs follow these definitions, rather than a change of basis or a
convention for boosts imported from another text. The numbering matches
the question sheet exactly.
\end{conventionbox}
\begingroup\small\tableofcontents\endgroup
\clearpage
\problem{1}{Deriving the Poincar\'e algebra from the group law}
\partsol{a}{Identity, inverse, and dimension}
Expanding \(\Lambda^T\eta\Lambda=\eta\) gives
\(\omega^T\eta+\eta\omega=0\), or \(\omega_{\mu\nu}=-\omega_{\nu\mu}\).
The identity is \((I,0)\), and direct composition gives
\[
 (\Lambda,a)^{-1}=(\Lambda^{-1},-\Lambda^{-1}a).
\]
There are \(D(D-1)/2\) independent Lorentz generators and \(D\)
translations, hence \(D(D+1)/2\) generators (ten for \(D=4\)).

\partsol{b}{Translations and the mixed bracket}
The translation subgroup is additive and abelian, so \([P^\mu,P^\nu]=0\).
The group law gives
\[
 S(\Lambda,0)S(I,a)S(\Lambda,0)^{-1}=S(I,\Lambda a).
\]
Because the exponent is \(-ia_\mu P^\mu\), comparison of coefficients
of the arbitrary covector \(a_\mu\) gives
\(S(\Lambda,0)P^\rho S(\Lambda,0)^{-1}
=(\Lambda^{-1})^\rho{}_{\sigma}P^\sigma\).
Expand this conjugation using the problem's convention for S. One obtains
\begin{equation}
 [M^{\mu\nu},P^\rho]
 =i(\eta^{\nu\rho}P^\mu-\eta^{\mu\rho}P^\nu).
 \label{sol4:MP}
\end{equation}
The inverse in the finite adjoint formula is consistent with the specified
order of conjugation; reversing that order gives the forward matrix.

\partsol{c}{Lorentz brackets}
Write \(X_j=-\frac i2\omega_{j\mu\nu}M^{\mu\nu}\).
The group commutator is
\((1+X_2)(1+X_1)(1-X_2)(1-X_1)
=1+[X_2,X_1]+O(\omega^3)\).
On coordinates the corresponding term is \([\omega_2,\omega_1]\).
Equating its independent coefficients yields
\begin{equation}
 [M^{\mu\nu},M^{\rho\sigma}]
 =i(\eta^{\nu\rho}M^{\mu\sigma}-\eta^{\mu\rho}M^{\nu\sigma}
       -\eta^{\nu\sigma}M^{\mu\rho}+\eta^{\mu\sigma}M^{\nu\rho}).
 \label{sol4:MM}
\end{equation}
This also follows by multiplying the supplied vector matrices: in
matrix-unit notation their entries are sums of \(E_{ab}\), and
\(E_{ab}E_{cd}=\delta_{bc}E_{ad}\) gives precisely the four terms above.
For example \([M^{12},M^{23}]=-iM^{13}\). Terms quadratic in only
\(\omega_1\) or only \(\omega_2\) cancel in the group commutator; the
mixed second-order term comes from their two first-order pieces.

\partsol{d}{Rotation and boost form}
With \(p_i=P^i\), the complete nonzero brackets are
\begin{align*}
 [J_i,J_j]&=i\epsilon_{ijk}J_k,&[J_i,K_j]&=i\epsilon_{ijk}K_k,
 &[K_i,K_j]&=-i\epsilon_{ijk}J_k,\\
 [J_i,p_j]&=i\epsilon_{ijk}p_k,&[K_i,H]&=ip_i,
 &[K_i,p_j]&=i\delta_{ij}H.
\end{align*}
Also \([J_i,H]=[H,p_i]=[p_i,p_j]=0\).
Thus \([K_1,K_2]=-iJ_3\). For \(K'_i=-K_i\), both brackets with
\(H,p_j\) reverse sign, while the boost--boost bracket and the form of
\([J_i,K'_j]\) do not. Reverse the boost parameter too to keep S unchanged.

\partsol{e}{Covering groups and projective phases}
The assumed exact multiplication law excludes a multiplier
\(e^{i\alpha(g_2,g_1)}\). A half-integer-spin representation is an ordinary
representation of the spin cover; the nontrivial covering element acts as
\(-1\). This global sign does not change the local Lie algebra or require
an added central generator. A projective representation of the quotient
and a Lie-algebra central extension are different issues.

\problem{2}{Euclidean rotations and translations}
\partsol{a}{Group law and generators}
Composition gives \((R_2,a_2)(R_1,a_1)=(R_2R_1,a_2+R_2a_1)\).
The specified pullback action has the same product law. Translations act
as \(e^{-ia^bP_b}f(x)=f(x-a)\). Choose a positive rotation in the
\(ab\) plane to send \(e_a\) toward \(e_b\). Then
\(U(R_{ab}(\theta),0)=e^{-i\theta L_{ab}}\): its first variation is
\(-\theta(x_a\partial_b-x_b\partial_a)f\), as required by \(R^{-1}x\).

\partsol{b}{Differential-operator brackets}
Using \([x_a,\partial_b]=-\delta_{ab}\), while derivatives commute,
\begin{align*}
 [P_a,P_b]&=0,\\
 [L_{ab},P_c]&=i(\delta_{ac}P_b-\delta_{bc}P_a),\\
 [L_{ab},L_{cd}]&=i(\delta_{ac}L_{bd}-\delta_{ad}L_{bc}
                       -\delta_{bc}L_{ad}+\delta_{bd}L_{ac}).
\end{align*}
For instance the mixed bracket follows by applying the product rule to
\(-i(x_a\partial_b-x_b\partial_a)(-i\partial_c)f\) and subtracting the
opposite order. The second derivatives cancel. The same product rule in
the rotation bracket cancels second-order differential operators and
leaves the four displayed first-order terms. There are \(n(n-1)/2\)
rotations and \(n\) translations. The latter form an abelian ideal,
because commutation with any generator returns a translation.

\partsol{c}{Checks and the missing translation sector}
The three requested tests are
\[
 [L_{12},P_1]=iP_2,\qquad [L_{12},P_2]=-iP_1,\qquad
 [L_{12},L_{23}]=-iL_{13}.
\]
Rotations fix the origin, whereas translations do not. No linear
combination of the rotation generators can supply a constant vector
field everywhere. The connected isometry group is \(ISO(n)\), not \(SO(n)\).

\partsol{d}{The Laplacian and action}
Let \(\Delta=\partial_c\partial_c\). Clearly \([P_a,\Delta]=0\), and
\[
 [x_a\partial_b-x_b\partial_a,\Delta]
 =-2\partial_a\partial_b+2\partial_b\partial_a=0.
\]
Under \(x'=Rx+a\), \(\phi'(x')=\phi(x)\), the gradient norm is invariant
by \(R^TR=I\) and the integration measure is unchanged. Hence \(S_E\)
is invariant. The same proof uses \(|\det R|=1\) and also covers
reflections for this scalar action, but a reflection is not generated
by the connected Lie algebra. Invariance of this particular action under
reflections must not be inferred from connected invariance alone.

\partsol{e}{Rotations versus boosts}
Writing \(B_i=L_{in}\) for the extra Euclidean plane gives
\([B_i,B_j]=+iL_{ij}\). With a timelike direction, the Lorentz metric
instead gives \([K_i,K_j]=-iM^{ij}\). Euclidean rotations have bounded,
periodic one-parameter subgroups; boosts have unbounded rapidity and
preserve an indefinite form. Dimension alone cannot identify the groups.

\problem{3}{The invariant norm of a vector in spinor notation}
\partsol{a}{Determinant, invariance, and inverse map}
Direct evaluation gives
\[
 \det V=(v^0+v^3)(v^0-v^3)-(v^1-iv^2)(v^1+iv^2)
       =(v^0)^2-|\bm v|^2.
\]
Since \(\det A=1\), \(\det(AVA^\dagger)=\det V\).
The trace identities of the Pauli matrices invert the map:
\(v^0=\tfrac12\tr V\), \(v^i=\tfrac12\tr(\sigma_iV)\).
Hermiticity makes these four numbers real.

\partsol{b}{Epsilon contraction and causal type}
Only the 12 and 21 components of each epsilon contribute. Expanding
all four terms gives
\[
 \tfrac12\varepsilon^{\alpha\beta}\varepsilon^{\dot\alpha\dot\beta}
 V_{\alpha\dot\alpha}V_{\beta\dot\beta}
 =\tfrac12(2V_{1\dot1}V_{2\dot2}-2V_{1\dot2}V_{2\dot1})=\det V.
\]
The two choices \(\varepsilon^{12}=-1\) multiply, so there is no
additional minus sign. Positive, negative, and zero determinant describe
timelike, spacelike, and null vectors, respectively (with the zero vector
separated if a nonzero null direction is meant). This is an invariant
quadratic form, not a positive Hilbert-space norm.

\partsol{c}{One spinor makes a null vector}
For \(\lambda=(a,b)^T\neq0\), \(\lambda\lambda^\dagger\) is positive
semidefinite of rank one, with determinant zero and positive trace.
Its vector is
\[
 v^0=\frac{|a|^2+|b|^2}{2},\quad v^1=\Re(a^*b),\quad
 v^2=\Im(a^*b),\quad v^3=\frac{|a|^2-|b|^2}{2}.
\]
It is future null. A timelike or spacelike vector cannot be represented
by this one outer product; nor can a past-directed nonzero null vector.

\partsol{d}{Two spinors distinguish causal types}
Let \(X=(\lambda\ \mu)\). Then
\(V_+=XX^\dagger\), \(V_-=X\diag(1,-1)X^\dagger\), and
\(\det X=\lambda^TE\mu\). Therefore
\[
 \det V_+=|\lambda^TE\mu|^2>0,\qquad
 \det V_-=-|\lambda^TE\mu|^2<0
\]
for independent columns. The first matrix is positive definite, so its
vector is future timelike. The second is spacelike. Dependent columns
instead give a zero determinant and include null or zero cases.

\partsol{e}{Tensor products and reality}
The first index transforms with A and the dotted index with its complex
conjugate, giving \((\tfrac12,0)\otimes(0,\tfrac12)=(\tfrac12,\tfrac12)\).
The tensor-product space consists of sums of decomposable tensors; a
single rank-one tensor is a special element, not a general one. The
complexified representation has four complex components. The condition
\(V^\dagger=V\) selects the four-real-dimensional vector representation.

\problem{4}{The rotation double cover and its Lorentz extension}
\partsol{a}{Orthogonality and orientation}
Unitary conjugation preserves Hermiticity and trace, so the image of
\(\bm v\cdot\bm\sigma\) again has real coefficients in the Pauli basis.
Since \(\tfrac12\tr[(\bm v\cdot\bm\sigma)^2]=|\bm v|^2\), it preserves
lengths and, by polarization, all dot products. Thus \(R(U)^TR(U)=I\).
The determinant is continuously valued in \(\{\pm1\}\), equals +1 at
U=I, and is +1 everywhere because \(SU(2)\) is connected.

\partsol{b}{Homomorphism and surjectivity}
Successive conjugations give \(R(U_2U_1)=R(U_2)R(U_1)\).
Use \(U=\cos(\theta/2)I-i\sin(\theta/2)\bm n\cdot\bm\sigma\)
and \((\bm a\cdot\bm\sigma)(\bm b\cdot\bm\sigma)
=\bm a\cdot\bm b+i(\bm a\times\bm b)\cdot\bm\sigma\). Multiplication yields
\[
 R(U)\bm v=\bm v\cos\theta+(\bm n\times\bm v)\sin\theta
       +\bm n(\bm n\cdot\bm v)(1-\cos\theta).
\]
This is Rodrigues' formula. Every proper three-dimensional rotation has
an axis and angle, so every such rotation is in the image.

\partsol{c}{Kernel and covering}
A matrix in the kernel commutes with all three Pauli matrices and hence
with the full matrix algebra. It is \(cI\); unitarity and determinant one
give \(c=\pm1\). The differential is an isomorphism between the
three-dimensional Lie algebras. A neighborhood of each of these two
kernel points maps diffeomorphically onto a neighborhood of the identity,
and translation gives the same local property elsewhere. Thus this is
a two-sheeted cover. At \(\theta=2\pi\), U is \(-I\) while R is I;
at \(4\pi\), U returns to I.

\partsol{d}{The Lorentz extension}
The map \(V\mapsto AVA^\dagger\) preserves Hermiticity and determinant,
so it preserves the Minkowski quadratic form. Continuity places its image
in the proper orthochronous component. In \(A=UH\), write
\(H=\exp(\rho\bm n\cdot\bm\sigma/2)\): positivity and determinant one
make its two eigenvalues reciprocal positive numbers. Conjugation by H
mixes \(v^0\) with \(\bm n\cdot\bm v\) by a hyperbolic rotation of rapidity
\(\rho\), while U gives the spatial rotation. Rotations and boosts generate
\(SO^+(1,3)\), proving surjectivity. If A acts trivially on every Hermitian
V, first use V=I to obtain unitarity, then the Pauli matrices to obtain
A=\(\pm I\). The kernel is again \(\mathbb Z_2\).

\partsol{e}{Which representations descend?}
A \(2\pi\) rotation acts in \((j_L,j_R)\) as
\((-1)^{2j_L+2j_R}\). Descent requires this to equal +1, equivalently
\(j_L+j_R\in\mathbb Z\). Scalars \((0,0)\) and vectors
\((\tfrac12,\tfrac12)\) descend; either Weyl representation does not.
The Dirac direct sum has \(-I\) on each block and also does not descend
as an ordinary representation of the Lorentz quotient.

\problem{5}{Sigma identities and the Clifford algebra}
\partsol{a}{Mixed identities}
For \(\mu=\nu=0\) the sum is \(2I\). For one time and one spatial index,
\(I(-\sigma_i)+\sigma_iI=0\). For two spatial indices,
\(-\sigma_i\sigma_j-\sigma_j\sigma_i=-2\delta_{ij}I\).
Reversing the order of the two sets gives the same answers. The mixed
products compose maps between the two Weyl spaces and return to one
space. An unmixed anticommutator would fail already for the time--space
case, where it would give \(2\sigma_i\).

\partsol{b}{Four-component consequences}
Multiplying the off-diagonal blocks of \(\gamma^\mu\) gives the two mixed
identities on its diagonal. Thus \(\{\gamma^\mu,\gamma^\nu\}=2\eta^{\mu\nu}I_4\).
Symmetric \(p_\mu p_\nu\) annihilates the antisymmetric gamma product, giving
\((\slashed p)^2=p^2I_4\). Similarly the derivatives commute, so
\[
 (i\slashed\partial-m)(i\slashed\partial+m)
 =-\gamma^\mu\gamma^\nu\partial_\mu\partial_\nu-m^2
 =-(\Box+m^2)I_4.
\]

\partsol{c}{The spin generator and vector bilinear}
Anticommuting one gamma through two others yields
\[
 [[\gamma^\mu,\gamma^\nu],\gamma^\rho]
 =4(\eta^{\nu\rho}\gamma^\mu-\eta^{\mu\rho}\gamma^\nu),\qquad
 [S^{\mu\nu},\gamma^\rho]
 =i(\eta^{\nu\rho}\gamma^\mu-\eta^{\mu\rho}\gamma^\nu).
\]
For \(\mathcal S=\exp[-\frac i2\omega_{\mu\nu}S^{\mu\nu}]\), this gives
\(\mathcal S^{-1}\gamma^\rho\mathcal S=\Lambda^\rho{}_{\sigma}\gamma^\sigma\).
Also \(\mathcal S^\dagger\gamma^0=\gamma^0\mathcal S^{-1}\).
Consequently \(\bar\Psi=\Psi^\dagger\gamma^0\) transforms as
\(\bar\Psi\mathcal S^{-1}\), and its gamma bilinear transforms as a vector.
Using \(\Psi^\dagger\) alone would not give the Lorentz adjoint for boosts.

\partsol{d}{Rotations and boosts in the two blocks}
Direct multiplication gives
\[
 S^{ij}=\frac12\epsilon_{ijk}\begin{pmatrix}\sigma_k&0\\0&\sigma_k\end{pmatrix},
 \qquad S^{i0}=\frac i2\begin{pmatrix}-\sigma_i&0\\0&\sigma_i\end{pmatrix}.
\]
Thus J is \(\bm\sigma/2\) in each block, whereas K has opposite imaginary
signs. The rotation generators are Hermitian, but these finite-dimensional
boost generators are not. This is consistent with the field representation
being nonunitary, not with a failure of unitarity of the particle Hilbert space.

\partsol{e}{Chirality and the mass term}
In this basis \(\gamma_*=\diag(-I,I)\). Hence \(\gamma_*^2=I\), it
anticommutes with each off-diagonal gamma, and it commutes with every
product of two gammas, including S. The projectors are
\(P_L=(I-\gamma_*)/2\), \(P_R=(I+\gamma_*)/2\).
Since \(\gamma^0\) anticommutes with \(\gamma_*\),
\(\overline{P_L\Psi}=\bar\Psi P_R\). Therefore
\[
 \bar\Psi\Psi=\overline{\Psi_L}\Psi_R+\overline{\Psi_R}\Psi_L;
 \qquad \overline{\Psi_L}\Psi_L=\overline{\Psi_R}\Psi_R=0.
\]
The scalar Dirac mass couples opposite chiralities.

\problem{6}{Pauli--Lubanski invariants and massless little groups}
\partsol{a}{Commuting invariants}
Translation commutativity gives \([P^2,P^\rho]=0\). Contracting
\eqref{sol4:MP} with \(P_\rho\) gives \([M^{\mu\nu},P^2]=0\), since the
antisymmetric combination \(P^\mu P^\nu-P^\nu P^\mu\) vanishes.
Likewise
\[
 [W^\mu,P^\lambda]=-\tfrac i2\epsilon^{\mu\nu\rho\sigma}P_\nu
 (\delta_\sigma{}^\lambda P_\rho-\delta_\rho{}^\lambda P_\sigma)=0,
 \qquad P_\mu W^\mu=0.
\]
In both equations a symmetric pair of momenta contracts with epsilon.
Applying \eqref{sol4:MP} and \eqref{sol4:MM} to the two factors in W,
the terms on the three contracted indices cancel by invariance of the
four-index epsilon tensor under infinitesimal proper Lorentz transformations.
The remaining free-index transformation is
\[
 [M^{\alpha\beta},W^\mu]
 =i(\eta^{\beta\mu}W^\alpha-\eta^{\alpha\mu}W^\beta).
\]
Contraction then gives \([M^{\alpha\beta},W_\mu W^\mu]=0\), even though
the W components need not commute: the two symmetric anticommutators
cancel. Since W also commutes with P, \(W^2\) is a Casimir.

\partsol{b}{The epsilon contraction and ordering}
When moving the second momentum past the first M in \(W^2\), the
commutator contributes a term of the form
\(\epsilon^{\mu\nu\rho\sigma}P_\nu
(\delta_\sigma{}^\lambda P_\rho-\delta_\rho{}^\lambda P_\sigma)\),
which vanishes before the second epsilon is contracted. Thus the momenta
may both be placed on the left in this particular expression. Now use
\[
 \epsilon^{\mu\nu\rho\sigma}\epsilon_{\mu\alpha\beta\gamma}
 =-\delta^{\nu\rho\sigma}_{\alpha\beta\gamma}.
\]
Expanding the six terms of the generalized Kronecker delta and using
antisymmetry of each M, the two terms pairing \(\nu\) with \(\alpha\)
give \(2P^2M^{\rho\sigma}M_{\rho\sigma}\); the remaining four give
\(-4P_\mu P_\nu M^{\mu\rho}M^\nu{}_{\rho}\).
Multiplication by \(-1/4\) gives
\[
 \boxed{W^2=-\tfrac12P^2M^{\mu\nu}M_{\mu\nu}
             +P_\mu P_\nu M^{\mu\rho}M^\nu{}_{\rho}.}
\]
No general license to commute P through M was used. The cancellation is
specific to the antisymmetric contractions.

\partsol{c}{Massive representations}
In the rest frame only \(P_0=m\) is nonzero, so \(W^0=0\) and
\(W^i=mJ_i\). Hence \(P^2=m^2\), \(W^2=-m^2\bm J^2=-m^2j(j+1)\).
These label mass and spin in a positive-energy massive irreducible
sector. They do not resolve extra internal multiplicities or, without
an energy-orientation assumption, the sign of energy.

\partsol{d}{The massless little algebra}
The brackets from Problem 1 give
\[
 [J,N_1]=iN_2,\quad [J,N_2]=-iN_1,\quad [N_1,N_2]=0.
\]
For example \([K_1,K_2]=-iJ_3\) cancels \([J_2,-J_1]=+iJ_3\)
in the last bracket. These are the \(\mathfrak{iso}(2)\) relations.
The general component formula is
\(W^0=\bm p\cdot\bm J\), \(\bm W=H\bm J+\bm p\times\bm K\).
At \(k=(\kappa,0,0,\kappa)\) it reduces to
\[
 W^\mu=\kappa(J,-N_2,N_1,J),\qquad
 W^2=-\kappa^2(N_1^2+N_2^2).
\]
These combinations stabilize k; for example the actions of K and the
corresponding transverse rotation on its components cancel.

\partsol{e}{Helicity and the finite-polarization assumption}
In a finite-dimensional unitary space the commuting Hermitian matrices
\(N_1,N_2\) can be simultaneously diagonalized. Conjugation by
\(e^{-i\theta J}\) rotates their joint eigenvalue pairs. A finite set of
points in a plane cannot contain the full circular orbit of a nonzero
point. Hence every pair is zero, and \(N_1=N_2=0\).
An irreducible polarization sector is then a character of rotations,
J=h, and \(W^\mu=hP^\mu\). All these helicities have \(P^2=W^2=0\);
the Casimir values do not determine h. For the spin cover, a \(4\pi\)
rotation is the identity and \(h\in\tfrac12\mathbb Z\).
Continuous-spin representations allow an infinite-dimensional polarization
space with nonzero translation orbit, and have \(W^2<0\).

\problem{7}{Weyl masses and a conserved Dirac charge}
\partsol{a}{Grassmann signs and Hermiticity}
Antisymmetry of E and the exchange of two odd components cancel:
\(B(\chi,\xi)=\chi_aE_{ab}\xi_b=-\xi_bE_{ab}\chi_a=B(\xi,\chi)\).
Thus \(B(\xi,\xi)=2\xi_1\xi_2\); for commuting spinors it would be zero
and the mixed bilinear would instead be antisymmetric.
Since \((s^\mu)^\dagger=s^\mu\),
\[
 (i\xi^\dagger s^\mu\partial_\mu\xi)^\dagger
 =-i(\partial_\mu\xi^\dagger)s^\mu\xi
 =i\xi^\dagger s^\mu\partial_\mu\xi
  -i\partial_\mu(\xi^\dagger s^\mu\xi).
\]
The action is Hermitian after the stated boundary term is dropped.

\partsol{b}{Mass phase and first-order equations}
Put \(\xi_{\rm old}=e^{-i\alpha/2}\xi_{\rm new}\). The two fields in
B supply \(e^{-i\alpha}\), so m becomes \(|m|\) while the kinetic term
is unchanged. A full-angle rotation would overcompensate the phase.
Before rephasing, \(B^\dagger=-2\xi_1^*\xi_2^*\), so
\(\Lag_m=-m\xi_1\xi_2+m^*\xi_1^*\xi_2^*\).
Putting each odd variation on the left gives
\[
 is^\mu\partial_\mu\xi+m^*E\xi^*=0.
\]
Varying the unconjugated variables after integration by parts gives its
conjugate. Define \(r=-E\xi^*\); using \(E(s^\mu)^*=\widetilde s^\mu E\)
the pair is
\[
 is^\mu\partial_\mu\xi-m^*r=0,\qquad
 i\widetilde s^\mu\partial_\mu r-m\xi=0.
\]
Acting on the first with \(i\widetilde s^\nu\partial_\nu\) and using the
second gives \((\Box+|m|^2)\xi=0\). The squared equation gives the
mass shell but loses the first-order relation between components and
the conjugate field; its arbitrary solutions are not all Weyl solutions.

\partsol{c}{Two equal masses and the flavor rotation}
The kinetic flavor metric and mass matrix are both proportional to the
identity, so a real orthogonal rotation of \((\xi_1,\xi_2)\) preserves
the free action. Choose
\(R(\alpha)=\left(\begin{smallmatrix}\cos\alpha&\sin\alpha\\
-\sin\alpha&\cos\alpha\end{smallmatrix}\right)\).
Then \(\chi\mapsto e^{-i\alpha}\chi\),
\(\zeta\mapsto e^{i\alpha}\zeta\), and
\[
 B(\xi_1,\xi_1)+B(\xi_2,\xi_2)=2B(\chi,\zeta).
\]
The kinetic term is \(i\chi^\dagger s\cdot\partial\chi+
 i\zeta^\dagger s\cdot\partial\zeta\); the mass term is
\(-m[B(\chi,\zeta)+B(\chi,\zeta)^\dagger]\).
The transformation between the two pairs of fields is invertible and
complex linear. Since \(\xi_1,\xi_2\) are independent complex Weyl fields,
\(\zeta\) is not constrained to be \(\chi^*\).

\partsol{d}{Dirac packaging and current}
For \(r=-E\zeta^*\), \(\bar\Psi=(r^\dagger,\chi^\dagger)\).
One finds \(\bar\Psi\Psi=B(\chi,\zeta)+B(\chi,\zeta)^\dagger\).
The lower kinetic block obeys
\[
 r^\dagger\widetilde s^\mu\partial_\mu r
 =\zeta^T(s^\mu)^T\partial_\mu\zeta^*
 =-(\partial_\mu\zeta^\dagger)s^\mu\zeta,
\]
where the last step exchanges odd variables. Integration by parts
therefore supplies the \(\zeta\) kinetic action. This proves
\(S=\int d^4x\,\bar\Psi(i\slashed\partial-m)\Psi\).
Both blocks of \(\Psi\) transform by \(e^{-i\alpha}\); Noether's current is
\[
 j^\mu=\bar\Psi\gamma^\mu\Psi
       =\chi^\dagger s^\mu\chi-\zeta^\dagger s^\mu\zeta.
\]
The minus sign follows from the same odd reordering without a derivative.
The Dirac equation and its adjoint give \(\partial_\mu j^\mu=0\).

\partsol{e}{When the charge is broken}
A term \(\mu B(\chi,\chi)+\hc\) acquires \(e^{-2i\alpha}\), so a
nonzero fixed \(\mu\) breaks this continuous phase symmetry, leaving
its sign subgroup. More generally \(\Xi\mapsto e^{-i\alpha Q}\Xi\)
changes a symmetric Weyl mass matrix to
\(e^{-i\alpha Q^T}Me^{-i\alpha Q}\). Its first variation vanishes exactly if
\[
 Q^TM+MQ=0.
\]
The kinetic term also requires Hermitian Q. For the Dirac pair,
\(Q=\diag(1,-1)\) and \(M=\left(\begin{smallmatrix}0&m\\m&0\end{smallmatrix}\right)\)
satisfy the condition; a nonzero 11 entry does not. Interactions must
separately respect the symmetry, so a degenerate free mass spectrum
alone is insufficient to conserve the charge in the full theory.

\problem{8}{Free spinors, spin sums, and the one-particle inner product}
\partsol{a}{Square roots and the Dirac equations}
Let \(X=\bm p\cdot\bm\sigma\), so \(X^2=\bm p^2I\).
Squaring the given numerators and using \(E^2-\bm p^2=m^2\) gives
\(A^2=EI+X\), \(B^2=EI-X\), and \(AB=BA=mI\).
Both matrices are Hermitian and positive for \(m>0\). In block form,
\[
 \slashed p=\begin{pmatrix}0&A^2\\B^2&0\end{pmatrix}.
\]
It sends \((A\xi,B\xi)^T\) to \(m(A\xi,B\xi)^T\), and
\((A\eta,-B\eta)^T\) to \(-m(A\eta,-B\eta)^T\).
This proves both required equations.

\partsol{b}{Three different orthogonality statements}
Matrix multiplication and orthonormality of each two-spinor basis give
\begin{align*}
 \bar u_su_t&=2m\delta_{st},&\bar v_sv_t&=-2m\delta_{st},
 &\bar u_sv_t&=0,\\
 u_s^\dagger u_t&=2E\delta_{st},&v_s^\dagger v_t&=2E\delta_{st},
 &u_s^\dagger(p)v_t(p)&=2\xi_s^\dagger X\eta_t.
\end{align*}
The last expression need not vanish. But \(A_{-p}=B_p\),
\(B_{-p}=A_p\), so \(u_s^\dagger(\bm p)v_t(-\bm p)=0\).
The vectors \(u(\bm p)\) and \(v(-\bm p)\) are eigenvectors of the same
Hermitian \(h_D(\bm p)\) with eigenvalues \(E\) and \(-E\). This is the
orthogonality implied by Hermiticity; \(v(\bm p)e^{-i\bm p\cdot\bm x}\)
has spatial Fourier momentum \(-\bm p\).

\partsol{c}{Spin sums and projectors}
Completeness \(\sum_s\xi_s\xi_s^\dagger=\sum_s\eta_s\eta_s^\dagger=I\)
gives
\[
 \sum_su_s\bar u_s=\begin{pmatrix}mI&A^2\\B^2&mI\end{pmatrix}
 =\slashed p+m,
 \quad
 \sum_sv_s\bar v_s=\begin{pmatrix}-mI&A^2\\B^2&-mI\end{pmatrix}
 =\slashed p-m.
\]
The slash-eigenvalue projectors are
\(\Pi_+=(m+\slashed p)/(2m)\), \(\Pi_-=(m-\slashed p)/(2m)\).
They sum to I and square to themselves by \(\slashed p^2=m^2\), but
are generally not Hermitian in the ordinary component inner product.
In particular the v sum must be divided by \(-2m\), not \(+2m\), to give
\(\Pi_-\). In contrast, the Hermitian energy projectors are
\[
 P_+(\bm p)=\frac{\sum_su_s(p)u_s^\dagger(p)}{2E}
   =\frac12\left(I+\frac{h_D(\bm p)}E\right),\qquad
 P_-(\bm p)=\frac{\sum_sv_s(p_-)v_s^\dagger(p_-)}{2E}
   =\frac12\left(I-\frac{h_D(\bm p)}E\right).
\]

\partsol{d}{The one-particle Hilbert space}
The invariant measure cancels the state normalization:
\[
 I_1=\sum_s\int d\Pi_p\,|\bm p,s\rangle\langle\bm p,s|,\quad
 |f\rangle=\sum_s\int d\Pi_p\,f_s(\bm p)|\bm p,s\rangle,\quad
 \|f\|^2=\sum_s\int d\Pi_p\,|f_s|^2.
\]
A normalized wave packet sets the last expression to one. This is a
positive norm of particle states with their invariant integration measure.
The invariant component bilinear \(\bar u u\) has a different role; its
v analogue is negative without implying negative-norm antiparticle states.

\Needspace{18\baselineskip}
\partsol{e}{The intertwiner}
Choose rest spinors \(u_s(k)\) and define
\(u_s(p)=\mathcal S(L(p))u_s(k)\). Insert the identity
\(L(\Lambda p)L(\Lambda p)^{-1}\) into \(\Lambda L(p)\), and write
\(W=L(\Lambda p)^{-1}\Lambda L(p)\). Since W fixes k, it is a rest-frame
rotation and acts on its spinors by the unitary matrix \(D^{(1/2)}(W)\).
Therefore
\[
 \mathcal S(\Lambda)u_s(p)
 =\sum_tu_t(\Lambda p)D^{(1/2)}_{ts}(W).
\]
Although \(\mathcal S\) on four components is nonunitary for boosts,
D is unitary and the mass-shell measure is invariant. The induced action
on f thus preserves the Hilbert-space norm in part (d). Covariantly
normalized states need no additional square-root energy factor.

\end{document}
